\begin{ex}
Tính: $345 + 234$.
\shortans{579}
\loigiai{
Ta có $345 + 234 = 579$.
}
\end{ex}

\begin{ex}
Tính: $678 - 345$.
\shortans{333}
\loigiai{
Ta có $678 - 345 = 333$.
}
\end{ex}

\begin{ex}
Tính: $456 + 278$.
\shortans{734}
\loigiai{
Ta có $456 + 278 = 734$.
}
\end{ex}

\begin{ex}
Tính: $734 - 456$.
\shortans{278}
\loigiai{
Ta có $734 - 456 = 278$.
}
\end{ex}

\begin{ex}
Tính: $528 + 365$.
\shortans{893}
\loigiai{
Ta có $528 + 365 = 893$.
}
\end{ex}

\begin{ex}
Tính: $820 - 456$.
\shortans{364}
\loigiai{
Ta có $820 - 456 = 364$.
}
\end{ex}

\begin{ex}
Tính: $600 - 245$.
\shortans{355}
\loigiai{
Ta có $600 - 245 = 355$.
}
\end{ex}

\begin{ex}
Tính: $1000 - 368$.
\shortans{632}
\loigiai{
Ta có $1000 - 368 = 632$.
}
\end{ex}

\begin{ex}
Tính: $1000 - 750$.
\shortans{250}
\loigiai{
Ta có $1000 - 750 = 250$.
}
\end{ex}

\begin{ex}
Tính giá trị biểu thức: $425 + 285 - 310$.
\shortans{400}
\loigiai{
Ta có $425 + 285 = 710$, rồi $710 - 310 = 400$.
}
\end{ex}

\begin{ex}
Tính giá trị biểu thức: $850 - 360 + 150$.
\shortans{640}
\loigiai{
Ta có $850 - 360 = 490$, rồi $490 + 150 = 640$.
}
\end{ex}

\begin{ex}
Tính giá trị biểu thức: $1000 - 450 - 250$.
\shortans{300}
\loigiai{
Ta có $1000 - 450 = 550$, rồi $550 - 250 = 300$.
}
\end{ex}

\begin{ex}
Điền số thích hợp: $350 + \dots = 800$.
\shortans{450}
\loigiai{
Số cần điền: $800 - 350 = 450$.
}
\end{ex}

\begin{ex}
Điền số thích hợp: $920 - \dots = 460$.
\shortans{460}
\loigiai{
Số cần điền: $920 - 460 = 460$.
}
\end{ex}

\begin{ex}
Điền số thích hợp: $\dots + 248 = 615$.
\shortans{367}
\loigiai{
Số cần điền: $615 - 248 = 367$.
}
\end{ex}

\begin{ex}
Điền số thích hợp: $\dots - 375 = 425$.
\shortans{800}
\loigiai{
Số cần điền: $425 + 375 = 800$.
}
\end{ex}

\begin{ex}
Trường Tiểu học Chu Văn An có 540 học sinh, trường Tiểu học Trần Phú có 430 học sinh. Cả hai trường có bao nhiêu học sinh?
\shortans{970}
\loigiai{
Cả hai trường có: $540 + 430 = 970$ học sinh.
}
\end{ex}

\begin{ex}
Kho A có 725 bao gạo, kho B có ít hơn kho A 150 bao gạo. Kho B có bao nhiêu bao gạo?
\shortans{575}
\loigiai{
Kho B có: $725 - 150 = 575$ bao gạo.
}
\end{ex}

\begin{ex}
Bác nông dân thu hoạch được 850 kg lúa, bác đã bán đi 560 kg. Bác còn lại bao nhiêu ki-lô-gam lúa?
\shortans{290}
\loigiai{
Bác còn lại: $850 - 560 = 290$ kg.
}
\end{ex}

\begin{ex}
Quãng đường từ Hà Nội đến Thanh Hóa dài khoảng 160 km, từ Thanh Hóa đến Vinh dài khoảng 140 km. Quãng đường từ Hà Nội đến Vinh đi qua Thanh Hóa dài bao nhiêu ki-lô-mét?
\shortans{300}
\loigiai{
Dài: $160 + 140 = 300$ km.
}
\end{ex}

\begin{ex}
Hai số có tổng bằng 950, số hạng thứ nhất là 420. Số hạng thứ hai là bao nhiêu?
\shortans{530}
\loigiai{
Số hạng thứ hai: $950 - 420 = 530$.
}
\end{ex}

\begin{ex}
Hai số có hiệu bằng 320, số trừ là 280. Số bị trừ là bao nhiêu?
\shortans{600}
\loigiai{
Số bị trừ: $320 + 280 = 600$.
}
\end{ex}

\begin{ex}
Tính tổng của số tròn trăm lớn nhất có ba chữ số và số 85.
\shortans{985}
\loigiai{
Số tròn trăm lớn nhất có 3 chữ số là 900. Tổng: $900 + 85 = 985$.
}
\end{ex}

\begin{ex}
Lấy 1000 trừ đi số tròn trăm liền sau 400 thì được bao nhiêu?
\shortans{500}
\loigiai{
Số tròn trăm liền sau 400 là 500. Hiệu: $1000 - 500 = 500$.
}
\end{ex}

\begin{ex}
Tính: $250 + 350 + 400$.
\shortans{1000}
\loigiai{
Ta có $250 + 350 = 600$, rồi $600 + 400 = 1000$.
}
\end{ex}
